\documentclass{ieeeaccess}
\usepackage{cite}
\usepackage{amsmath,amssymb,amsfonts}
\usepackage{algorithmic}
\usepackage{graphicx}
\usepackage{textcomp}
\usepackage{booktabs}
\usepackage{multirow}
\usepackage{url}
\usepackage{xcolor}
\def\BibTeX{{\rm B\kern-.05em{\sc i\kern-.025em b}\kern-.08em
    T\kern-.1667em\lower.7ex\hbox{E}\kern-.125emX}}
\def\xfigwd{\columnwidth}
\usepackage{stfloats}
\setcounter{topnumber}{4}
\setcounter{bottomnumber}{4}
\setcounter{totalnumber}{10}
\setcounter{dbltopnumber}{4}
\renewcommand{\topfraction}{0.95}
\renewcommand{\bottomfraction}{0.95}
\renewcommand{\textfraction}{0.05}
\renewcommand{\floatpagefraction}{0.9}
\renewcommand{\dbltopfraction}{0.95}
\renewcommand{\dblfloatpagefraction}{0.9}

\begin{document}
\history{Date of publication xxxx 00, 0000, date of current version xxxx 00, 0000.}
\doi{10.1109/ACCESS.2026.DOI}

\title{STORM-PhysNet: A Multi-Horizon Transformer for Geostationary Relativistic Electron Flux Forecasting with Physics-Inspired Components and Cross-Satellite Transfer}

\author{\uppercase{Samarth BN}\authorrefmark{1},
\IEEEmembership{Student Member, IEEE}}
\address[1]{RV College of Engineering, Bengaluru, India (e-mail: samarthbn.ec25@rvce.edu.in)}
\tfootnote{The authors thank the providers of GOES, OMNI, and GRASP data products.}
\tfootnote{This work extends a concurrent conference submission on the same system. The present manuscript adds the wider-delay ablation, the bagged-Transformer control, the architecture-matched Transformer control, expanded interpretability diagnostics, and the full multi-seed GRASP transfer analysis.}
\markboth
{Samarth BN: STORM-PhysNet for GEO Electron Flux Forecasting}
{Samarth BN: STORM-PhysNet for GEO Electron Flux Forecasting}
\corresp{Corresponding author: Samarth BN (e-mail: samarthbn.ec25@rvce.edu.in).}

\begin{abstract}
Energetic electrons above 2\,MeV at geostationary Earth orbit (GEO)
are a persistent internal-charging hazard for satellites. We present
STORM-PhysNet, a multi-horizon forecast model that couples a standard
Transformer encoder over a 72-hour window with a learnable L1--Earth
propagation delay, a \(B_z\)-conditioned physics gate, and residual
heads for simultaneous 1-hour, 6-hour, and 12-hour log-flux
forecasts.

Relative to an architecture-matched Transformer control
($d_{\mathrm{model}}=128$, two layers, four heads), STORM-Bz retains a small
1\,h PE$_{\mathrm{clim}}$ edge (0.986 vs 0.980) and a large advantage versus
persistence (PE$_{\mathrm{pers}} \approx 0.31$ vs $\approx -0.01$). At 6\,h and 12\,h, single-model
means are close (STORM-Bz 0.900 / 0.854 vs matched 0.895 / 0.845).
Seed-bagged STORM-Bz reaches PE$_{6\mathrm{h}} = 0.910$ and PE$_{12\mathrm{h}} = 0.870$;
bagged matched Transformer is competitive (0.908 / 0.861).
A default Transformer baseline ($d_{\mathrm{model}}=64$, three layers; not
capacity-matched) scores PE$_{1\mathrm{h}} = 0.978$ and is reported only as a
secondary reference. Controlled ablations (No-Delay, No-Gate,
No-Physics) show that the short-horizon gain is a property of the
trained package as a whole rather than of any individual physics
module at inference.

Under synthetic solar-wind input noise and targeted channel dropout,
STORM degrades more slowly than the Transformer, especially at
1-hour. On GSAT-19 GRASP at Indian longitude, zero-shot transfer is
weak at long horizons; fine-tuning on GRASP raises 6\,h PE from 0.740 to 0.841 (paired $p = 3.55 \times 10^{-12}$) and 12\,h PE from 0.567 to 0.762 ($p = 1.66 \times 10^{-13}$) on a $\sim$14--15 month GRASP window (fifteen seeds; paired tests on seed-level PE). We report bootstrap confidence intervals, paired tests, case studies, and release evaluation scripts so that later work can reuse the same splits and PE definitions.
\end{abstract}

\begin{keywords}
Geostationary orbit, GRASP, physics-inspired neural network, prediction
efficiency, radiation belt, relativistic electrons, space weather,
Transformer encoder, transfer learning
\end{keywords}

\titlepgskip=-15pt
\maketitle

\section{Introduction}
\label{sec:introduction}
\PARstart{T}{he} outer electron radiation belt at geostationary Earth orbit
(GEO) is one of the most dynamically variable plasma environments in near-Earth
space. Relativistic electrons with energies \(E>2\,\mathrm{MeV}\) can undergo
flux variations spanning four to five orders of magnitude over timescales of
hours to days, driven by the interplay of solar wind forcing, magnetospheric
wave-particle interactions, and global current systems~\cite{b1,b2,b3}. These
populations represent a chronic engineering hazard: energetic electrons
penetrating satellite shielding accumulate charge on internal dielectrics and
conductors, potentially triggering electrostatic discharges that damage or
destroy spacecraft electronics~\cite{b1}. With the rapid expansion of the
operational GEO belt---hosting communication, broadcasting, and navigation
assets---accurate forecasting of electron fluxes has become an operational
priority for satellite operators and space agencies worldwide.

Forecast models for GEO electrons have evolved considerably over four decades.
The earliest approaches used linear autoregressive filters and empirical
correlations with solar wind speed and geomagnetic indices~\cite{b4,b5}.
NOAA's Relativistic Electron Forecast Model (REFM), which has operated since
the 1990s, remains a reference benchmark for the community~\cite{b5}.
We do not include REFM as a numeric baseline because it targets daily integral
fluence rather than hourly log-flux; a direct PE comparison would be
misleading~\cite{b5}.
The subsequent generation of feed-forward neural networks demonstrated that
nonlinear mappings from solar wind and IMF parameters could improve upon these
linear baselines~\cite{b7,b8}. LSTM-based architectures then raised the bar
further by explicitly capturing multi-day temporal dependencies in the
solar wind driving~\cite{b9,b10,b11,b12}. More recently, Transformer and
attention-based models---originally developed for natural language
processing~\cite{b13}---have been adapted for space weather time series,
reporting daily-fluence prediction efficiencies of 0.85--0.92 with or without
pre-training strategies~\cite{b14,b15,b16}.

Despite this rapid progress, several gaps remain. First, most published models
target either a single forecast horizon (typically 24\,h fluence) or a small
number of multi-hour lead times; the short-horizon regime (about 1\,h) that
is critical for immediate operational response has received less attention as a
stand-alone multi-seed evaluation target. Second, the community lacks a
controlled multi-seed benchmark that directly compares physics-inspired and
physics-agnostic architectures at equivalent training budgets on the same
chronological splits. Third, nearly all published work focuses on GOES satellites
at American/European longitudes; the performance of transfer strategies for
Indian-longitude GEO assets such as ISRO's GSAT-19---which carries the
Geostationary Radiation Alert and Monitoring System (GRASP) instrument---has not
been systematically characterized.

STORM-PhysNet is designed around these three gaps.

\subsection{Contributions}
\begin{itemize}
\item A Transformer encoder with adaptive L1--Earth delay and a
\(B_z\)-conditioned gate, evaluated at 1\,h, 6\,h, and 12\,h under
fifteen matched seeds and chronological splits.
\item Evidence that STORM-Bz retains a small 1\,h PE$_{\mathrm{clim}}$ edge
versus an architecture-matched Transformer (0.986 vs 0.980) and a large
PE$_{\mathrm{pers}}$ gap. Validation selects a mixing weight \(\alpha^*\) per seed (typically 0.4--0.6). We also report a fixed \(\alpha=0.3\) diagnostic. Under selected \(\alpha^*\), linear-ensemble test PE\(_{6\mathrm{h}} \approx 0.906\). We do not claim online \(\alpha\) adaptation.
\item Seed-bagged STORM as a pure-STORM multi-horizon alternative that
approaches ensemble skill without a second architecture at inference.
\item Noise and missing-channel robustness comparisons, GRASP
zero-shot versus fine-tune transfer with fifteen seeds, and a null
ablation showing that short-horizon-only physics losses do not move PE.
\item Case studies, delay/gate diagnostics, and released evaluation
scripts for reuse of the same PE definitions and splits.
\end{itemize}


\section{Related Work}
\label{sec:related}

\subsection{Physics of Outer Radiation Belt Electrons at GEO}
The outer radiation belt at GEO (\(L\approx6.6\)) is maintained by a dynamic
balance between electron acceleration and loss. Storm-time acceleration and loss
are reviewed in detail by Reeves et al.~\cite{b17} and Thorne~\cite{b18}. Which
storms produce large GEO electron responses is discussed by O'Brien et
al.~\cite{b19}; broader source--loss modeling is reviewed by Shprits et
al.~\cite{b20}. Acceleration occurs primarily through whistler-mode chorus wave
interactions near the magnetic equator~\cite{b2,b3,b6}. Ultra-low frequency (ULF)
waves additionally drive inward radial diffusion. Loss processes include
precipitation into the atmosphere and magnetopause shadowing during storm
compressions. The net response of GEO flux to a geomagnetic storm depends on the
relative balance of these mechanisms and the direction and magnitude of the
interplanetary magnetic field \(B_z\)~\cite{b2}. This motivates our choice to
condition the primary gate module on \(B_z\).

\subsection{Machine Learning for GEO Electron Flux}
Camporeale's review framed both the opportunities and the methodological pitfalls
of data-driven space weather modeling~\cite{b21}. Morley et al.\ defined and
advocated for prediction efficiency as the preferred scalar skill score for
relativistic electron forecasts~\cite{b22}.

Early neural network models established that solar wind speed, IMF \(B_z\), and
geomagnetic indices carry predictive signal for GEO electrons~\cite{b7}. Wei et
al.\ demonstrated that a deep network on multi-day solar wind histories achieves
strong daily PE~\cite{b8}. Son et al.\ pushed to 72-hour multi-step hourly
forecasts with LSTM~\cite{b9}. Zhang et al.\ combined empirical mode decomposition
with LSTM~\cite{b10}. Simms et al.\ focused on classification of enhancement
events~\cite{b11}.

Transformer-based architectures entered the GEO flux literature around
2023--2024~\cite{b14,b15,b16}. Stacking ensembles have also been explored for
short-horizon energetic electron prediction~\cite{b23}. The iTransformer of Liu
et al.~\cite{b24} was briefly evaluated as an alternative backbone; under the same
multi-seed protocol it showed higher variance across seeds than the standard
temporal-attention Transformer, which was therefore retained for all reported
experiments.

Beyond GEO, Sinha et al.\ developed PreMevE~\cite{b25}; Botek
et al.~\cite{b26} applied neural networks at low Earth orbit; Wang et
al.~\cite{b12} used LSTM with MagEIS observations.

A recurring theme in this literature is that headline PE depends strongly on the
chosen cadence, energy channel, solar-cycle phase, and persistence definition.
Daily fluence scores near 0.9 are not directly comparable to hourly log-flux PE
on a different year range~\cite{b21,b22}. Our evaluation therefore emphasizes
internal controls: the same chronological splits, the same fifteen seeds, and the
same PE formula for every architecture, including ablations. That design does not
replace a community benchmark, but it does make relative statements among the
models in this paper interpretable.

Physics-inspired machine learning for radiation-belt forecasting is still less
common than pure sequence models. Where physics appears, it is often encoded as
post-hoc diagnostics, hard-wired lag features, or loss penalties rather than as
trainable modules whose outputs can be inspected (for example, a learned transit
delay whose distribution can be compared to the L1--Earth travel-time window).
STORM-PhysNet is deliberately conservative on the encoder side---a standard
temporal Transformer---so that interpretability claims attach to the delay and
gate, not to an exotic backbone.

\begin{table*}[!htbp]
\centering
\caption{Design comparison with closely related GEO electron flux forecasting studies.
PE numbers are omitted on purpose: they are not interchangeable across different
periods, cadences, and baselines~\cite{b21,b22}.}
\label{tab:compare}
\begin{tabular}{@{}lp{2.0cm}p{2.3cm}p{4.0cm}@{}}
\toprule
\textbf{Study} & \textbf{Core model} & \textbf{Horizon focus} & \textbf{Relative to this work} \\
\midrule
Wei et al.\ (2018) & LSTM & 1-day fluence & Daily target; no multi-seed protocol \\
Son et al.\ (2022) & MLP & 1--72\,h hourly & No physics modules; no cross-satellite transfer \\
Zhang et al.\ (2022) & EMD-LSTM & 1-day & Decomposition pipeline; daily resolution \\
Tan et al.\ (2024) & TPA-LSTM / TF & Up to 3\,days (5-min) & Multi-energy channels; no delay/gate ablations \\
Sun et al.\ (2024) & LSTM / TF & 1\,h--1\,day & Closest multi-horizon TF baseline; no module ablations or GRASP transfer \\
Sun et al.\ (2025) & TF + data strategies & Daily fluence & Fluence target; few-shot/pretrain focus \\
Sinha et al.\ (2021) & PreMevE ensemble & $\sim$1--2\,day belt-wide & Multi-mission belt model; not GEO-only multi-horizon heads \\
\textbf{This work} & TF + delay + \(B_z\) gate & 1\,h / 6\,h / 12\,h & 15-seed protocol; honest module ablations; GOES$\rightarrow$GRASP transfer \\
\bottomrule
\end{tabular}
\end{table*}

\subsection{Solar Wind Propagation Delay}
Solar-wind transit from L1 to the magnetopause is typically on the order of 0.5--1.5\,h~\cite{b27} depending on solar-wind speed. The delay module is constrained to \([0.5, 1.5]\,\mathrm{h}\), a band chosen to cover that range while still allowing limited storm-time deviation without breaking causality. We do not claim a unique recovered constant of 1.09\,h from every window.

\subsection{Transfer Learning and Cross-Satellite Adaptation}
Most radiation-belt transfer efforts have focused on cross-calibration of flux
measurements across fleets~\cite{b25}. The specific problem of adapting a model
trained on GOES-15 to GSAT-19 GRASP in a different longitude sector has not been
previously reported in the form studied here. We explore transfer by adapting
a GOES-pretrained model to GRASP with limited local labels.

\section{Data and Problem Setup}
\label{sec:data}

\subsection{GOES Satellite Data}
We use measurements of \(E>2\,\mathrm{MeV}\) integral electron flux from the
GOES-15 satellite. The chronological working sample spans approximately March 2010
to February 2020. Raw one-minute flux values are aggregated to hourly means;
the within-hour standard deviation is retained as an input feature so that
intra-hour variability is not discarded. All flux values are log-transformed
(\(\log_{10}\) of particle cm\(^{-2}\)s\(^{-1}\)sr\(^{-1}\)).

\subsection{OMNI Solar-Wind and Geomagnetic Inputs}
Solar wind and interplanetary magnetic field parameters are drawn from the OMNI
1-hour dataset. The released dataloader uses sixteen solar-wind / geomagnetic channels plus log-flux history (17 inputs total after concatenation):
\begin{itemize}
\item Plasma and pressure: \(V_{sw}\), proton density, \(P_{\mathrm{dyn}}\)
\item Magnetic field: \(B_t\), \(B_z\)
\item Derived activity: \(B_z\)-south duration, storm-onset hours, \(B_z\times V_{sw}\)
\item Rolling moments: \(V_{sw}\) at 1/6/24\,h, \(B_z\) at 1/6\,h
\item Indices: Kp, Dst
\item Within-hour variability: \(\mathrm{std}(\log\mathrm{flux})\) over the 1-minute samples
\end{itemize}
Forecast targets use integer hourly leads of 1\,h, 6\,h, and 12\,h on the hourly GOES series (config/dataloader: $[1,6,12]$). A 72-hour history window (\(T=72\)) is used at each forecast origin.

\begin{table}[h]
\centering
\caption{Input feature inventory used by the released dataloader}
\label{tab:features}
\begin{tabular}{@{}ll@{}}
\toprule
\textbf{Category} & \textbf{Channels} \\ \midrule
Plasma \& pressure & \(V_{sw}\), density, \(P_{\mathrm{dyn}}\) \\
Magnetic field & \(B_t\), \(B_z\) \\
Derived drivers & \(B_z\)-south duration, storm-onset hours, \(B_z\times V_{sw}\) \\
Rolling moments & \(V_{sw}\) (1/6/24\,h), \(B_z\) (1/6\,h) \\
Geomagnetic indices & Kp, Dst \\
Within-hour flux variability & \(\log\)-flux std over 1-min samples \\
Autoregressive target history & log-flux (separate input stream) \\ \bottomrule
\end{tabular}
\end{table}

\subsection{Chronological Splits}
Splits are purely chronological: the earliest 70\% of the merged GOES--OMNI series
for training, the next 15\% for validation, and the final 15\% for held-out
testing. No shuffling is applied.

\subsection{Forecast Targets and Horizons}
The model simultaneously predicts log-flux at three lead times:
\begin{equation}
\hat{\mathbf{y}}_{t}
=
\bigl[\hat{y}_{t+1\mathrm{h}},\,
\hat{y}_{t+6\mathrm{h}},\,
\hat{y}_{t+12\mathrm{h}}\bigr].
\label{eq:horizons}
\end{equation}
The 1\,h head targets an operationally actionable warning window.
The 6\,h head matches the lead time most commonly benchmarked in the literature.
The 12\,h head tests medium-range predictability.

\subsection{Evaluation Metrics}
We report two PE definitions so that results remain comparable both to
climatological skill and to persistence-based operational baselines
following Morley et al.~\cite{b22} and the methodological cautions in
Camporeale~\cite{b21}:
\begin{align}
\mathrm{PE}_{\mathrm{clim}}
&= 1 - \frac{\mathrm{MSE}(\hat y,y)}{\mathrm{Var}(y)}, \\
\mathrm{PE}_{\mathrm{pers}}
&= 1 - \frac{\mathrm{MSE}(\hat y,y)}{\mathrm{MSE}(y^{\mathrm{pers}},y)}.
\end{align}
Unless stated otherwise, headline GOES tables use \(\mathrm{PE}_{\mathrm{clim}}\)
with fifteen seeds and bootstrap 95\% confidence intervals. The fifteen seeds measure sensitivity to random initialization on a fixed chronological split; they do not constitute independent trials across different solar-cycle phases or data periods. \(\mathrm{PE}_{\mathrm{pers}}\) is emphasized at the 1-hour horizon,
where persistence is a strong operational baseline.
Storm windows use Dst \(\le -50\,\mathrm{nT}\).

\subsection{GRASP Dataset}
GRASP aboard GSAT-19 provides GEO radiation measurements at Indian longitude
(\(\approx\)82\(^\circ\)E). We use the available record covering approximately
July 2017 to September 2018---an $\approx$14--15 month interval in the declining phase of Solar
Cycle 24. After alignment, GRASP supplies fewer co-located input channels than the
GOES pipeline; this feature mismatch is a proximate cause of zero-shot failure and
motivates input-projection adaptation. Storm windows on the GRASP test set are
defined by Dst \(\le -50\,\mathrm{nT}\) (118 of 1781 test hours, 6.6\%).

\section{Method: STORM-PhysNet}
\label{sec:method}

\subsection{Overall Architecture}
\begin{figure*}[!t]
\centering
\includegraphics[width=0.95\textwidth]{figures/fig_system_architecture.png}
\caption{STORM-PhysNet: OMNI solar wind is delay-aligned ($\tau\in[0.5,1.5]\,\mathrm{h}$)
and concatenated with GOES log-flux ($T=72\,\mathrm{h}$), projected to
$d_{\mathrm{model}}=128$, and encoded by a 2-layer, 4-head Transformer.
A $B_z$-conditioned gate modulates the final representation; three residual heads
with shared horizon embeddings ($d_e=16$) predict corrections over persistence
at 1\,h, 6\,h, and 12\,h.}
\label{fig:arch}
\end{figure*}

STORM-PhysNet follows a delay--encode--gate--decode design
(Fig.~\ref{fig:arch}). For a batch of input sequences
\(\mathbf{X}\in\mathbb{R}^{B\times T\times F}\):
(1) the adaptive delay module interpolates the solar-wind channels by \(\tau\)
time steps;
(2) an input projection maps the concatenated solar-wind and flux features at each
time step to dimension \(d_{\mathrm{model}}\);
(3) a Transformer encoder processes the sequence and the final time-step
representation is extracted;
(4) the \(B_z\)-conditioned gate computes a per-channel modulation vector;
(5) three independent residual heads produce the three horizon forecasts.
Measured parameter counts are approximately $3.43\times10^{5}$ for
STORM-Bz and $8.45\times10^{5}$ for the Transformer baseline
(default hyperparameters). The two models differ in both width and depth;
they are not architecture-matched. We additionally trained an architecture-matched
Transformer control with the same encoder size as STORM
($d_{\mathrm{model}}=128$, two layers, four heads;
$\approx1.19\times10^{6}$ parameters) under the same chronological split and
fifteen seeds. Default-baseline PE differences should still be read with the
parameter-count disparity in mind; matched-control PE is reported separately.






\section{Results}
\label{sec:results}

Unless noted, all GOES numbers below are PE$_{\mathrm{clim}}$ on the fixed
chronological test split, aggregated over fifteen seeds (42--56).
Test PE is computed once after training and is never used for model selection.

\begin{table*}[!t]
\centering
\caption{GOES test PE$_{\mathrm{clim}}$ (seed mean). PE$_{\mathrm{pers},1\mathrm{h}}$ is
skill versus persistence; PE$_{\mathrm{st},6\mathrm{h}}$ is restricted to storm-flagged
hours. Bagged rows average predictions across fifteen seeds (true bagging).
Bold: best single-model or bagged entry in each column among reported systems.}
\label{tab:main}
\begin{tabular}{lccccc}
\toprule
System & PE$_{1\mathrm{h}}$ & PE$_{6\mathrm{h}}$ & PE$_{12\mathrm{h}}$ & PE$_{\mathrm{st},6\mathrm{h}}$ & PE$_{\mathrm{pers},1\mathrm{h}}$ \\
\midrule
LSTM & 0.955 & 0.881 & 0.840 & 0.788 & $-1.19$ \\
Transformer (default) & 0.978 & 0.895 & 0.845 & 0.797 & $-0.08$ \\
Transformer matched & 0.980 & 0.895 & 0.845 & 0.809 & $-0.01$ \\
STORM-Bz & \textbf{0.986} & 0.900 & 0.854 & 0.812 & \textbf{+0.31} \\
No-Delay & 0.987 & 0.902 & 0.855 & 0.809 & +0.36 \\
No-Physics & 0.986 & 0.899 & 0.850 & 0.807 & +0.31 \\
No-Gate & 0.986 & 0.900 & 0.856 & 0.819 & +0.32 \\
\midrule
STORM-Bz bagged & 0.987 & \textbf{0.910} & \textbf{0.870} & 0.836 & --- \\
TF matched bagged & 0.984 & 0.908 & 0.861 & 0.831 & --- \\
\bottomrule
\end{tabular}
\end{table*}

Table~\ref{tab:main} is the primary GOES leaderboard for this paper. The primary comparison in this section is between STORM-Bz and the architecture-matched Transformer control.
The largest single-model separation versus the default Transformer is at 1\,h
(PE$_{\mathrm{clim}}$ and especially PE$_{\mathrm{pers}}$). At 6--12\,h, seed means
are close; multi-seed aggregation (bagging) provides the clearer multi-horizon lift.

\begin{figure*}[!t]
\centering
\includegraphics[width=0.92\textwidth]{figures/fig_horizon_pe.png}
\caption{Per-horizon PE$_{\mathrm{clim}}$ (seed mean $\pm$ std) for the main systems
in Table~\ref{tab:main}. Short-horizon separation is visible; 6\,h/12\,h means overlap
more strongly across ablations.}
\label{fig:horizon_pe}
\end{figure*}

\subsection{Ablations and Matched Control}

In all STORM variants, physics-loss terms are restricted to the
short-horizon head by default, as early testing showed that enforcing
these constraints on the 6\,h and 12\,h heads did not improve long-horizon PE.

We further trained a No-Gate variant that sets the modulation
to identity (\(g=1\)). On PE\(_{\mathrm{clim}}\), removing the gate
does not reduce 1\,h skill relative to full STORM-Bz
(\(0.986\) versus \(0.986\) in Table~\ref{tab:main}); the same
null result holds for the No-Delay (which reaches \(0.987\)) and No-Physics
ablations. The short-horizon gain versus the Transformer
is therefore a property of the trained STORM package under our
protocol, not a gain that disappears when any one physics module is
removed. We do not attribute that 1\,h PE\(_{\mathrm{clim}}\) edge
to the gate alone.

On storm windows, single-model PE\(_{\mathrm{st},6\mathrm{h}}\) values are close:
STORM-Bz \(0.812\), LSTM \(0.788\), Transformer \(0.797\), with
No-Delay and No-Physics slightly lower
(\(0.823\) and \(0.824\)) and No-Gate comparable or slightly
higher (\(0.829\)). Given overlapping intervals, we do not claim that
any one physics module is required for storm-time skill. The clearest
gains on storm subsets remain post-hoc aggregation: the
ensemble under selected \(\alpha^*\) reaches \(0.822\) and seed-bagged STORM
reaches \(0.836\) (Table~\ref{tab:main}).

Fig.~\ref{fig:ablation} summarizes the same controlled comparison at 6\,h:
ablating delay or physics loss does not overturn the ranking already reported
in Table~\ref{tab:main}.

\begin{figure}[!htbp]
\centering
\includegraphics[width=\columnwidth]{figures/fig_ablation_6h.png}
\caption{Six-hour PE against climatology for the Transformer baseline,
full STORM-Bz, and ablations that remove the delay module or the physics-loss
term (fifteen seeds; mean\(\pm\)std). Removing the delay module, the physics-loss term, or the \(B_z\) gate
does not collapse PE\(_{6\mathrm{h}}\) or remove the 1\,h advantage
over the Transformer; the main single-model STORM advantage remains
the short-horizon comparison in Table~\ref{tab:main}.}
\label{fig:ablation}
\end{figure}

\begin{figure}[!htbp]
\centering
\includegraphics[width=\columnwidth]{figures/fig_grasp_domain_gap.png}
\caption{GRASP domain-gap summary (fifteen seeds): zero-shot GOES models versus
fine-tuning on GSAT-19. Short-horizon PE transfers more readily; 6\,h and 12\,h
PE improve substantially after adaptation.}
\label{fig:grasp}
\end{figure}





\begin{table}[!t]
\centering
\caption{Trainable parameter counts for the released configurations.}
\label{tab:params}
\begin{tabular}{lr}
\toprule
Model & Parameters \\
\midrule
LSTM & $5.47\times10^{4}$ \\
STORM-Bz / No-Delay / No-Gate / No-Physics & $3.43\times10^{5}$ \\
RDG (cathode--anode gate) & $3.41\times10^{5}$ \\
RDG-S (RDG + spectral head) & $3.51\times10^{5}$ \\
Transformer (default) & $8.45\times10^{5}$ \\
Transformer matched ($d{=}128$, 2 layers) & $1.19\times10^{6}$ \\
\bottomrule
\end{tabular}
\end{table}

Table~\ref{tab:params} shows that the default Transformer is wider/deeper than STORM. Default Transformer is not capacity-matched to STORM; matched control exceeds STORM in encoder size by design.
in total parameter count, while the matched control is deliberately larger than STORM
in encoder capacity. We therefore never interpret default-Transformer gaps as pure
evidence for the delay/gate modules alone.


\subsection{Alternative Gate Formulations}
\label{sec:alt_gates}

Beyond the default $B_z$ gate we evaluated three probes under the same
fifteen-seed protocol:
\begin{itemize}
  \item \textbf{RDG} (rectified driver gate): cathode--anode style modulation driven by
        southward driver intensity;
  \item \textbf{RDG-S}: RDG with an auxiliary spectral residual head;
  \item \textbf{SDG} (saturating dose gate): radiotrophic-style saturating nonlinearity.
\end{itemize}
Seed-mean PE values remain within noise of STORM-Bz at all three horizons
(Table~\ref{tab:alt_gates}). We treat these variants as sensitivity checks on the
gating nonlinearity, not as replacements for the default $B_z$ gate in the main claim.

\begin{table}[!t]
\centering
\caption{Alternative gates (seed-mean PE$_{\mathrm{clim}}$, fifteen seeds).}
\label{tab:alt_gates}
\begin{tabular}{lccc}
\toprule
Gate & PE$_{1\mathrm{h}}$ & PE$_{6\mathrm{h}}$ & PE$_{12\mathrm{h}}$ \\
\midrule
STORM-Bz (default) & 0.986 & 0.900 & 0.854 \\
RDG & 0.986 & 0.901 & 0.855 \\
RDG-S & 0.986 & 0.902 & 0.856 \\
SDG & 0.986 & 0.901 & 0.854 \\
\bottomrule
\end{tabular}
\end{table}

\begin{figure}[!t]
\centering
\includegraphics[width=\columnwidth]{figures/fig_alt_gates_pe6h.png}
\caption{Alternative gates at 6\,h (seed mean $\pm$ std). Differences relative to STORM-Bz are small compared with seed scatter.}
\label{fig:alt_gates}
\end{figure}

\subsection{Interpretability}
\begin{figure*}[!t]
\centering
\includegraphics[width=0.95\textwidth]{figures/fig_physics_tau_hist.png}
\caption{Fraction of predicted L1--Earth delays \(\tau\) on the GOES test set across fifteen seeds. Blue: interior of \([0.5,1.5]\,\mathrm{h}\); red: near the upper bound (\(1.5\,\mathrm{h}\)); green: near the lower bound. Saturation at \(1.5\,\mathrm{h}\) shows the box constraint is active; this is not evidence of a unique recovered transit time.}
\label{fig:delay}
\end{figure*}

\begin{figure}[!htbp]
\centering
\includegraphics[width=\columnwidth]{figures/fig_physics_gate_storm_quiet.png}
\caption{\(B_z\) gate activation: storm windows (orange) versus quiet intervals
(blue). The absolute activation band is narrow ($\sim$0.38--0.45), showing weak $B_z$ dependence. This is consistent with the null gate ablation which showed no large PE drop when the gate was removed.}
\label{fig:gate}
\end{figure}

Fig.~\ref{fig:physics_scatter} relates the two trainable physics modules to
familiar solar-wind drivers; these panels are diagnostic, not a claim that
\(\tau\) recovers a unique physical constant on every window.

\begin{figure*}[!t]
\centering
\includegraphics[width=0.95\textwidth]{figures/fig_physics_gate_activation.png}
\caption{Physics-module diagnostics on the GOES test set: predicted delay
\(\tau\) against solar-wind speed, and gate activation against southward
\(B_z\) statistics. Delays remain inside the imposed \([0.5,1.5]\,\mathrm{h}\)
band; gate values rise when southward IMF is sustained.}
\label{fig:physics_scatter}
\end{figure*}

\begin{figure}[!htbp]
\centering
\includegraphics[width=\columnwidth]{figures/fig_feature_importance.png}
\caption{Permutation importance on the \(6\,\mathrm{h}\) head for STORM-Bz
(PE drop when each input channel is shuffled). \(B_z\) and related solar-wind
drivers rank highest.}
\label{fig:perm}
\end{figure}



\subsection{Case Study and Seed Stability}

\begin{figure*}[!t]
\centering
\includegraphics[width=0.95\textwidth]{figures/fig_case_study_timeseries.png}
\caption{Six-hour forecasts during a major geomagnetic-storm flux enhancement in the
GOES test window (STORM-Bz versus architecture-matched Transformer). The model tracks
the multi-day rise and recovery more closely than persistence; we do not claim that
gate activations alone identify reconnection intervals from this plot.}
\label{fig:case_enhancement}
\end{figure*}

\begin{figure}[!t]
\centering
\includegraphics[width=\columnwidth]{figures/fig_seed_spread.png}
\caption{Seed-to-seed PE spread for the main systems. Tight clusters indicate that
headline means are not driven by a single lucky initialization.}
\label{fig:seed_spread}
\end{figure}

\begin{figure}[!t]
\centering
\includegraphics[width=\columnwidth]{figures/fig_dst_bins.png}
\caption{PE stratified by Dst regime (diagnostic on the working GOES--OMNI sample).
Storm bins remain harder than quiet bins for all systems; absolute values are
illustrative and are not used for model selection.}
\label{fig:dst_bins}
\end{figure}

\subsection{GRASP Transfer}
On GSAT-19 GRASP (fifteen seeds), zero-shot GOES models already achieve
short-horizon PE near \(0.96\), but 6\,h and 12\,h PE fall to
\(0.740\) and \(0.567\). Fine-tuning on GRASP (freezing the Transformer
encoder and updating only the forecast heads) raises 6\,h PE from 0.740
to 0.841 (paired \textit{t}-test across fifteen seeds,
\(p = 3.55\times10^{-12}\), Cohen's \(d = 5.60\))
and 12\,h PE from 0.567 to 0.762
(\(p = 1.66\times10^{-13}\), \(d = 7.00\)).
The inter-seed standard deviation after fine-tuning decreases to 0.017 at 6\,h
and 0.020 at 12\,h, indicating that transfer converges more stably than zero-shot.
Absolute GRASP PE remains below
the GOES test scores, as expected under sensor and longitude shift; the
practical point is the large mid- and long-horizon recovery after
adaptation rather than a claim of parity with GOES.

\begin{table*}[!htbp]
\centering
\caption{GRASP (GSAT-19) transfer results (fifteen seeds). Paired tests across fifteen seeds are reported in the text. Bold: best per horizon.}
\label{tab:grasp}
\begin{tabular}{lccc}
\toprule
Horizon & Zero-shot & Fine-tuned & Gain \\
\midrule
1 h  & 0.964 \(\pm\) 0.001 & 0.968 \(\pm\) 0.001 & +0.005 \\
6 h  & 0.740 \(\pm\) 0.020 & \textbf{0.841 \(\pm\) 0.017} & \textbf{+0.102} \\
12 h & 0.567 \(\pm\) 0.033 & \textbf{0.762 \(\pm\) 0.020} & \textbf{+0.195} \\
\bottomrule
\end{tabular}
\end{table*}


\section{Discussion}
\label{sec:discussion}
The controlled comparison supports a layered reading rather than a single
``winner'' architecture. At 1\,h, STORM-Bz improves reliability relative to the architecture-matched Transformer, especially against persistence. Controlled ablations that remove the delay, the physics loss, or the gate do not by themselves erase that short-horizon PE\(_{\mathrm{clim}}\) edge. We therefore treat the gain as a property of the full trained system rather than as proof that any single module is necessary.

As detailed in Section~\ref{sec:results}, the ablations together indicate that the short-horizon gain is produced by the overall training protocol rather than by any single runtime module. The delay and gate act as structured inductive biases during training; once the model is trained, removing either module at inference does not erase the 1\,h PE\(_{\mathrm{clim}}\) edge over the Transformer. STORM-PhysNet should therefore be understood primarily as a carefully engineered training and multi-horizon recipe that includes interpretable physics-inspired components. To test whether the original upper bound of \(1.5\,\mathrm{h}\) was limiting performance, we retrained STORM with wider delay constraints (\(2.0, 2.5, 3.0, 3.5,\) and \(4.0\,\mathrm{h}\)) across fifteen seeds (sixteen seeds for the \(2.0\,\mathrm{h}\) bound). Mean PE$_{1\mathrm{h}}$ remained in the narrow range \(0.9859\)--\(0.9862\) and PE$_{6\mathrm{h}}$ remained in the range \(0.900\)--\(0.902\) for all bounds (see Fig.~\ref{fig:wider_delay}). The original \([0.5, 1.5]\,\mathrm{h}\) constraint is therefore not a performance bottleneck.

\begin{figure*}[!t]
\centering
\includegraphics[width=0.95\textwidth]{figures/wider_delay_pe6h.png}
\caption{Wider delay ablation: mean PE$_{6\mathrm{h}}$ versus delay upper bound
(error bars show seed standard deviation). Vertical dashed lines mark the
paper baselines STORM-Bz (\(0.900\)) and Transformer (\(0.895\)).
PE$_{6\mathrm{h}}$ remains stable across upper bounds from \(1.5\,\mathrm{h}\) to \(4.0\,\mathrm{h}\),
confirming that the original \([0.5, 1.5]\,\mathrm{h}\) constraint is not a performance bottleneck.}
\label{fig:wider_delay}
\end{figure*}

An architecture-matched Transformer control (\(d_{\mathrm{model}}=128\), two layers)
reached mean PE$_{1\mathrm{h}}=0.980$ and PE$_{6\mathrm{h}}=0.895$ across
fifteen seeds; bagging the fifteen matched seeds yields PE$_{1\mathrm{h}}=0.984$,
PE$_{6\mathrm{h}}=0.908$, and PE$_{12\mathrm{h}}=0.861$. Relative to STORM-Bz
(\(0.986\) / \(0.900\) / \(0.854\)), matching encoder width/depth shrinks the short-horizon
gap and shifts the multi-horizon advantage toward the larger Transformer under
bagging. This reinforces that single-model PE gaps are protocol-plus-width/depth
effects rather than proof that the delay or gate modules are individually necessary.
The package-level claim is falsifiable in this study: an architecture-matched Transformer, three module ablations, and true bagging are all reported under the same split and seeds.

At 6--12\,h, a plain Transformer remains competitive as a
single model; combining the two predictors (or bagging STORM seeds)
provides the clearest multi-horizon gain among the systems evaluated here.
Seed-bagged STORM reaches a multi-horizon PE level comparable to the linear ensemble (test PE\(_{6\mathrm{h}} \approx 0.906\) under selected \(\alpha^*\)).
Noise and missing-channel probes indicate that the physics-inspired model
fails more gracefully when upstream solar-wind features are corrupted.
GRASP results show that Indian-longitude deployment still needs local
adaptation for longer leads, but short-horizon skill transfers more
readily than long-horizon skill.

Dst intensity bins reinforce that reading. Skill does not vanish in moderate or
strong storm bins; the strong-storm bin is small (\(n=206\)), so uncertainty is
higher, but PE remains clearly above zero. Robustness grids show a smooth
degradation under noise rather than a brittle cliff at the first perturbation,
which is desirable for a research prototype even though it is not a substitute
for hardware-in-the-loop testing.

\subsection{Transfer Learning for Indian Longitude}
The GRASP experiments illustrate a practical constraint for new GEO instruments:
a GOES-pretrained model can retain useful short-horizon skill zero-shot, but
mid- and long-horizon PE drop until local fine-tuning is applied.
In our fifteen-seed runs, fine-tuning raises 6\,h PE from \(0.740\) to \(0.841\) and 12\,h PE from \(0.567\) to \(0.762\). Absolute GRASP PE
remains below the GOES test scores, as expected under sensor and longitude shift. A similar adaptation profile is observed for the architecture-matched Transformer baseline, which improves its 6\,h PE from \(0.797 \pm 0.014\) (zero-shot) to \(0.878 \pm 0.006\) after fine-tuning.

\subsection{Comparability With Published PE Numbers}
PE values reported across studies should not be compared without normalization
for data period, solar-cycle phase, energy channel, cadence, and persistence
baseline~\cite{b21,b22}. Daily fluence scores near 0.9 on one archive are not
evidence that an hourly log-flux model on a 2010--2020 sample is ``behind.'' Our training
period spans a moderate activity phase of Solar Cycle 24; results during a more
active phase may differ. The contribution of this paper is a controlled
multi-seed comparison and a documented transfer recipe.

\section{Limitations}
\label{sec:limits}
All reported leads are integer hours on hourly data; sub-hour (e.g.\ 45\,min) supervision would require higher-cadence flux labels.
Main GOES tables use fifteen seeds with bootstrap intervals; some
secondary diagnostics (individual case windows, older single-seed
robustness grids) remain thinner. PE values are not interchangeable with
daily-fluence scores reported elsewhere without matching cadence, energy
channel, and baseline~\cite{b21,b22}. GRASP covers a limited declining-phase
window. Performance during solar maximum (Cycle 25, 2023--2025) is unknown and likely to differ, particularly in the storm-time bins. Ensemble skill uses a validation-selected \(\alpha^*\) (with \(\alpha=0.3\) as a fixed diagnostic); the test-set \(\alpha\) curve is diagnostic only.
We do not claim online \(\alpha\) adaptation. Horizon-restricted
physics losses did not improve PE in our tests. The reported p-values are uncorrected for multiple comparisons across horizons and ablations; treat individual significance claims as descriptive rather than confirmatory.
The default Transformer baseline is not architecture-matched to STORM; the matched
control (\(d_{\mathrm{model}}=128\), two layers) shows that encoder width/depth and bagging
account for much of the multi-horizon PE movement, so we still do not attribute
single-model gaps solely to the delay or gate modules.

Alternative gates (RDG, RDG-S, SDG) remain within seed noise of STORM-Bz and are reported as sensitivity checks, not primary systems.

\section{Conclusion}
\label{sec:conclusion}
STORM-PhysNet is a reproducible multi-horizon baseline on GOES with a documented
transfer path toward GSAT-19 GRASP. Relative to the architecture-matched Transformer, the clearest single-model gain is short-horizon
reliability (PE$_{1\mathrm{h}}$ and PE$_{\mathrm{pers}}$).
Among systems evaluated here, seed-bagged STORM-Bz (and the RDG/SDG probes)
reaches the strongest multi-horizon PE
(PE$_{6\mathrm{h}}=0.910$, PE$_{12\mathrm{h}}=0.870$).
An architecture-matched bagged Transformer remains competitive
(\(0.908\) / \(0.861\)) but does not dominate under this full retrain.
Input-noise and missing-channel checks, together with GRASP fine-tuning
gains at 6--12\,h, support operational motivation beyond a pure leaderboard PE
number; for instance, the architecture-matched baseline remains near persistence at 1\,h (PE$_{\mathrm{pers}}\approx-0.01$), whereas STORM-Bz is clearly positive (PE$_{\mathrm{pers}}\approx+0.31$). Furthermore, as a noise probe, under synthetic Gaussian noise injection (\(\sigma=1.0\)), the 6\,h PE of the matched baseline drops to \(0.825 \pm 0.011\), highlighting the need for structurally robust architectures during degraded sensor conditions. Horizon-restricted physics losses did not improve PE in a fifteen-seed
ablation.
\section*{Data Availability}
The code, configuration files, trained checkpoints, evaluation scripts,
result CSVs, paper figures, and the chronological split protocol are
available at
\url{https://github.com/bnsama29-cloud/STORM-PhysNet-L2}. Note that the GRASP fine-tuning run, wider-delay retrain, and robustness/interpretability diagnostics were executed externally, and their frozen outputs are provided in the repository.
The GOES-15 EPEAD dataset is publicly available from NOAA NCEI
(\url{https://www.ngdc.noaa.gov/stp/satellite/goes/}).
OMNI is available via NASA OMNIWeb
(\url{https://omniweb.gsfc.nasa.gov/}).
GSAT-19 GRASP is accessible via ISSDC
(\url{https://www.issdc.gov.in/}).
\section*{Acknowledgment}
The authors thank the data providers for GOES, OMNI, and GRASP products.

\begin{thebibliography}{99}
\bibitem{b1}
D.~N. Baker et al., ``Space weather effects in the Earth's radiation belts,''
\emph{Space Sci.\ Rev.}, vol.~214, no.~1, p.~17, 2018.

\bibitem{b2}
R.~B. Horne et al., ``Wave acceleration of electrons in the Van Allen
radiation belts,'' \emph{Nature}, vol.~437, pp.~227--230, 2005.

\bibitem{b3}
D.~L. Turner et al., ``Explaining sudden losses of outer radiation belt
electrons during geomagnetic storms,'' \emph{Nature Phys.}, vol.~8,
pp.~208--212, 2012.

\bibitem{b4}
D.~N. Baker et al., ``Linear prediction filter analysis of relativistic
electron properties at 6.6 \(R_E\),'' \emph{J. Geophys. Res.}, vol.~95,
pp.~15133--15140, 1990.

\bibitem{b5}
K.~A. Perry et al., ``Comparing geosynchronous relativistic electron
prediction models,'' \emph{Space Weather}, vol.~8, S12003, 2010.

\bibitem{b6}
J.~Bortnik and R.~M. Thorne, ``The dual role of ELF/VLF chorus waves in
the acceleration and precipitation of radiation belt electrons,''
\emph{J. Atmos. Sol.-Terr. Phys.}, vol.~69, pp.~378--386, 2007.

\bibitem{b7}
A.~G. Ling et al., ``Identifying solar wind drivers of radiation belt
responses,'' \emph{Space Weather}, vol.~8, no.~10, 2010.

\bibitem{b8}
L.~Wei et al., ``Quantitative prediction of high-energy electron integral
flux at geostationary orbit based on deep learning,'' \emph{Space Weather},
vol.~16, pp.~903--916, 2018.

\bibitem{b9}
J.~Son et al., ``72-hour time series forecasting of hourly relativistic
electron fluxes at geostationary orbit by deep learning,''
\emph{Space Weather}, vol.~20, e2022SW003153, 2022.

\bibitem{b10}
H.~Zhang et al., ``A prediction model of relativistic electrons at
geostationary orbit using the EMD-LSTM network,'' \emph{Space Weather},
vol.~20, e2022SW003126, 2022.

\bibitem{b11}
L.~E. Simms et al., ``Classifier neural network models predict relativistic
electron events at geosynchronous orbit,'' \emph{J. Geophys. Res.\
Space Physics}, vol.~125, e2019JA027323, 2020.

\bibitem{b12}
Y.~Wang et al., ``LSTM-based modeling of radiation belt electron dynamics
using MagEIS Van Allen Probes observations,'' \emph{Space Weather}, vol.~23,
e2024SW004001, 2025.

\bibitem{b13}
A.~Vaswani et al., ``Attention is all you need,'' in \emph{Proc.\ NeurIPS},
2017, pp.~5998--6008.

\bibitem{b14}
W.~Tan et al., ``Comparative analysis of TPA-LSTM and Transformer models
for forecasting GEO radiation belt electron fluxes,'' \emph{Space Weather},
vol.~22, e2023SW003714, 2024.

\bibitem{b15}
X.~Sun et al., ``A modeling study of \(\ge2\) MeV electron fluxes in GEO
at different prediction time scales based on LSTM and transformer networks,''
\emph{J. Space Weather Space Clim.}, vol.~14, p.~20, 2024.

\bibitem{b16}
X.~Sun et al., ``Enhancing \(\ge2\) MeV electron fluence predictions in GEO
through three deep learning strategies,'' \emph{Space Weather}, vol.~23,
e2024SW003975, 2025.

\bibitem{b17}
G.~D. Reeves et al., ``Acceleration and loss of relativistic electrons
during geomagnetic storms,'' \emph{Geophys.\ Res.\ Lett.}, vol.~30,
no.~10, 2003.

\bibitem{b18}
R.~M. Thorne, ``Radiation belt dynamics: The importance of wave-particle
interactions,'' \emph{Geophys.\ Res.\ Lett.}, vol.~37, L22107, 2010.

\bibitem{b19}
T.~P. O'Brien et al., ``Which magnetic storms produce relativistic
electrons at geosynchronous orbit?'' \emph{J. Geophys.\ Res.},
vol.~106, pp.~15533--15544, 2001.

\bibitem{b20}
Y.~Y. Shprits et al., ``Review of modeling of losses and sources of
relativistic electrons in the outer radiation belt,'' \emph{J. Atmos.
Sol.-Terr.\ Phys.}, vol.~70, pp.~1694--1713, 2008.

\bibitem{b21}
E.~Camporeale, ``The challenge of machine learning in space weather:
Nowcasting and forecasting,'' \emph{Space Weather}, vol.~17, pp.~1166--1207,
2019.

\bibitem{b22}
S.~K. Morley et al., ``Evaluating the performance of radiation belt models:
Metrics, baselines, and statistics,'' \emph{Space Weather}, vol.~16,
pp.~1521--1548, 2018.

\bibitem{b23}
R.~Tang et al., ``Short-time prediction of energetic electron flux based
on stacking ensemble learning,'' \emph{Space Weather}, vol.~20,
e2021SW002928, 2022.

\bibitem{b24}
Y.~Liu et al., ``iTransformer: Inverted Transformers are effective for time
series forecasting,'' in \emph{Proc.\ ICLR 2024} (Spotlight), 2024.
arXiv:2310.06625.

\bibitem{b25}
S.~Sinha et al., ``PreMevE update: Forecasting ultra-relativistic electrons
inside Earth's outer radiation belt,'' \emph{Space Weather}, vol.~19,
e2021SW002773, 2021.


\bibitem{b26}
E.~Botek et al., ``Prediction of radiation belt electron fluxes at LEO using
neural networks with PROBA-V/EPT data,'' \emph{Space Weather}, vol.~21,
e2023SW003466, 2023.

\bibitem{b27}
N.~Y.~Ganushkina et al., ``Forecasting GOES 15 \(>2\) MeV electron fluxes from
solar wind data and geomagnetic indices,'' \emph{Space Weather}, vol.~18,
e2019SW002416, 2020.

\bibitem{b28}
D.~P. Kingma and J.~Ba, ``Adam: A method for stochastic optimization,''
in \emph{Proc.\ Int. Conf. Learn. Represent.\ (ICLR)}, 2015.




\end{thebibliography}

\vspace{3ex}
\noindent\textbf{Samarth BN} is pursuing the B.E.\ degree in Electronics and
Communication Engineering at RV College of Engineering, Bengaluru, India.
His interests include space-weather forecasting, time-series machine learning,
and radiation-belt physics.

\EOD
\end{document}